% =========================================================================
% Chapter 7 - Physics-informed operators
% =========================================================================
\chapter{Physics-informed operators}
\label{ch:pifno}

\section{Constrain by construction, not by penalty}

A soft penalty $\lambda\lVert\text{violation}\rVert^2$ trades accuracy against
constraint satisfaction and fully achieves neither: the weight must be tuned,
the constraint holds only approximately, and a badly scaled term degrades the
model while looking principled. Where a constraint can be built into the
\emph{parameterisation} it should be. It then holds for every weight
configuration, at initialisation and after any optimiser step, with nothing to
tune and no inference cost.

Three constraints move into the architecture.

\subsection{Positivity}

Predict $\log\rho$ and exponentiate. The inverse transform returns
$e^{y}-\epsilon$, so $\rho > -\epsilon$ \emph{identically}.

\subsection{Electron count}

\begin{equation}
  \rho_{\mathrm{out}}(\rr) = N_{\mathrm{val}}\,
    \frac{\tilde\rho(\rr)}{\int \tilde\rho\,\mathrm{d}\rr}
\end{equation}
makes $\int\rho = N_{\mathrm{val}}$ hold to machine precision, differentiably,
at the cost of one reduction. Strictly better than a penalty on the same
quantity.

\subsection{The von Weizsäcker bound}
\label{sec:pauli-head}

Using the exact decomposition~\eqref{eq:pauli}, write the network output as
\begin{equation}
  \boxed{\;
    \tau_\theta = \tauvW[\rho] \;+\; s\cdot\mathrm{softplus}\bigl(f_\theta(\rho)\bigr) \;}
  \label{eq:paulihead}
\end{equation}
so the network learns only the Pauli term. This buys two things at once:

\begin{enumerate}
  \item \textbf{The bound becomes structural.} \opt{softplus} is non-negative,
    so $\tau_\theta \ge \tauvW$ identically. The inequality is enforced
    \emph{non-strictly}, exactly as Hoffmann-Ostenhof states it: for a strongly
    negative output \opt{softplus} underflows and the head returns
    $\tau=\tauvW$ exactly --- the correct one-orbital limit, which is
    \emph{reachable} rather than merely approachable. A soft penalty can only
    approach it from above.
  \item \textbf{The network stops re-deriving known physics.} On the gold data
    $\tauvW$ accounts for $\sim31\,\%$ of $\tau$; that fraction is now supplied
    analytically and exactly by a spectral gradient, leaving only the genuinely
    unknown remainder to learn.
\end{enumerate}

\begin{ptip}
The scale $s$ is fitted on the \emph{training split only} --- fitting it on the
held-out material would leak the answer --- and is then optimised as $\log s$,
which cannot change sign. $\tauvW[\rho]$ is a fixed function of the
\emph{input}, so it contributes no gradient to $\theta$: it is a per-sample
offset, not a second trainable branch.
\end{ptip}

\begin{pwarn}
The bound is a theorem for all-electron densities; \file{CHGCAR} and
\file{TAUCAR} are pseudo quantities, so it must be \emph{measured} before the
head is enabled on a new data set. \opt{pauli\_bound\_violation} does this.
On the present data it holds at every point of two structures and at all but
one point of the third.
\end{pwarn}

\section{Measured effect of the head}
\label{sec:headresult}

Identical data, seed, architecture and epoch count; only the head differs.

\begin{center}
\begin{tabular}{@{}lrrr@{}}
\toprule
Held out & baseline & \textbf{Pauli head} & change \\
\midrule
struct\_000 & $0.0534$ & $\mathbf{0.0363}$ & $-32.0\,\%$ \\
struct\_001 & $0.0840$ & $\mathbf{0.0762}$ & $-9.3\,\%$ \\
struct\_002 & $0.0902$ & $\mathbf{0.0734}$ & $-18.6\,\%$ \\
\midrule
mean & $0.0759$ & $\mathbf{0.0620}$ & $\mathbf{-18.3\,\%}$ \\
\bottomrule
\end{tabular}
\end{center}

(relative $L^2$ on the held-out material.) The head wins on every fold and every
metric; MSE falls $30.6\,\%$.

Two results deserve more attention than the headline number.

\begin{description}
  \item[$\Ts$ improves more than $\tau$ does.] The integrated kinetic energy ---
    the quantity that actually enters the total energy --- improves from
    $2.55\to1.13\,\%$, $0.53\to0.27\,\%$ and $0.90\to0.58\,\%$: roughly halved
    on every fold, a larger relative gain than the pointwise error. That is the
    expected direction, since $\tauvW$ is supplied exactly and contributes no
    model error to the integral.
  \item[It trains \emph{faster}.] $148.6$\,s per fold against $180.8$\,s, and to
    a lower final training loss. The extra FFT for $\tauvW$ is more than repaid
    by conditioning: the network fits only $\tauP$, whose dynamic range is far
    narrower than $\tau$'s. This is not an accuracy-versus-cost trade.
\end{description}

\begin{pwarn}
The constraint did \emph{not} bind. The unconstrained baseline violated the
bound at zero points in all three folds, with margins of $0.83$--$1.05\
\mathrm{eV}/\text{\AA}^3$. On converged reference data the head's gain therefore
comes from handing the model $31\,\%$ of the answer analytically, not from
rescuing violations. The hard guarantee is insurance that costs nothing here and
becomes load-bearing inside an SCF loop, where the functional sees
non-converged densities far from the training distribution. These three
structures cannot distinguish ``the head enforces the bound'' from ``the bound
is easy here''; the enforcement claim rests on the tests, which check it at
random initialisation, at $50\times$ weights and after training at a
destructive learning rate.
\end{pwarn}

\section{Soft constraints}

Where a constraint cannot be built in, it enters the objective:
\begin{equation}
  \mathcal{L} = \mathcal{L}_{\mathrm{data}}
    + \lambda_{\mathrm{EL}}\mathcal{L}_{\mathrm{EL}}
    + \lambda_{\mathrm{TF}}\mathcal{L}_{\mathrm{TF}}
    + \lambda_{T}\mathcal{L}_{T}
    + \lambda_{\mathrm{scale}}\mathcal{L}_{\mathrm{scale}} .
\end{equation}

Every physics weight defaults to zero, so the objective is exactly the
supervised baseline until a term is enabled deliberately. Three design rules:

\begin{enumerate}
  \item \textbf{Physics terms act on the prediction decoded to physical units},
    never on the normalised representation --- a constraint is a statement about
    physics, not about preprocessing.
  \item \textbf{Every term is dimensionless.} A raw
    $(\mathrm{eV}/\text{\AA}^3)^2$ penalty varies by orders of magnitude between
    a light semiconductor and a transition-metal oxide, and no single weight
    could serve both.
  \item \textbf{The data term is $H^1$, not $L^2$}, when gradients matter:
    $\tauvW$ depends on $\nabla\rho$, and gradient noise is amplified by
    $|\nabla\rho|^2/8\rho$ exactly where $\rho$ is small.
\end{enumerate}

\subsection{The Euler--Lagrange residual}

The most valuable soft term is Eq.~\eqref{eq:euler} itself. Since $\mu$ is a
constant, subtracting the cell average eliminates it and leaves a residual that
must vanish for the exact density:
\begin{equation}
  r(\rr) = \frac{\delta \Ts}{\delta\rho} + \vext + \vH[\rho] + \vxc[\rho]
           - \overline{(\cdots)} .
  \label{eq:elresidual}
\end{equation}
It contains only $\vext$ (the network input) and $\rho$ (its output), so it
requires \textbf{no additional labels}.

\begin{pwarn}
The residual is only as good as the kinetic functional used to build it. With
the analytic $\mathrm{TF}+\lambda\,\mathrm{vW}$ surrogate it supplies a
physically correct \emph{inductive bias}, not ground truth, and must carry a
modest weight.
\end{pwarn}

\section{The closed loop}

The end state is not two models but one differentiable orbital-free cycle. The
link is Eq.~\eqref{eq:euler} with $\delta\Ts/\delta\rho$ obtained by autograd
through Model~2, since $\Ts=\int\tau_{\theta_B}[\rho]\,\mathrm{d}\rr$:

\begin{center}
\begin{tikzpicture}[
    node distance=8mm and 12mm,
    box/.style={draw=poraquered, rounded corners=2pt, fill=poraquered!5,
                inner sep=5pt, font=\small},
    file/.style={draw=shadegray, rounded corners=2pt, fill=panelgray,
                 inner sep=5pt, font=\small\ttfamily},
    arr/.style={-{Latex[length=2mm]}, thick, draw=poraquedark}]
  \node[file] (ext) {EXTCAR};
  \node[box, right=of ext] (m1) {FNO 1: $\vext \mapsto \rho$};
  \node[file, right=of m1] (chg) {CHGCAR};
  \node[box, right=of chg] (m2) {FNO 2: $\rho \mapsto \tau$};
  \node[file, right=of m2] (tau) {TAUCAR};
  \draw[arr] (ext) -- (m1);
  \draw[arr] (m1) -- (chg);
  \draw[arr] (chg) -- (m2);
  \draw[arr] (m2) -- (tau);
  \node[below=14mm of chg, font=\small\itshape, text=shadegray]
    (loop) {Euler--Lagrange residual (label-free)};
  \draw[arr, draw=poraqueamber] (chg) |- ([yshift=-6mm]chg.south) -| (loop);
  \draw[arr, draw=poraqueamber, dashed] (m2.south) |- (loop.east);
  \draw[arr, draw=poraqueamber, dashed] (loop.west) -| (m1.south);
\end{tikzpicture}
\end{center}

Three consequences, in order of practical value:

\begin{enumerate}
  \item \textbf{Model~2 is regularised on Model~1's output distribution.} A
    functional trained only on converged reference densities fails when placed
    inside an SCF loop, because it then sees non-converged densities it was
    never shown. Joint training on the residual exposes it to exactly that
    distribution.
  \item \textbf{The optimised quantity becomes the used quantity.} Orbital-free
    DFT needs $\delta\Ts/\delta\rho$, not $\tau$; training through the
    derivative removes the objective mismatch of Chapter~\ref{ch:ofdft}.
  \item \textbf{Inference becomes physically consistent.} A predicted
    $(\rho,\tau)$ pair satisfies the variational condition it is supposed to
    satisfy, rather than merely resembling the training data.
\end{enumerate}

Beyond that, replacing Model~1's forward pass with a fixed-point solve of
Eq.~\eqref{eq:euler} using the learned functional --- a deep-equilibrium model,
differentiated by the implicit function theorem --- would make the network an
orbital-free \emph{solver} rather than a regressor, with transferability bounded
by the functional alone.

\section{Why this is cheap here}

Physics-informed training is normally expensive: derivatives must come from
autograd through the network or from finite differences, and the latter injects
a truncation error that the loss misattributes to the model. On a plane-wave
grid neither problem arises.

\begin{center}
\begin{tabular}{@{}lll@{}}
\toprule
Operator & On this grid & Cost \\
\midrule
$\nabla$ & $i\GG$ & one FFT, \textbf{exact} \\
$\nabla^2$ & $-G^2$ & one FFT, \textbf{exact} \\
$\vH[\rho]$ & $4\pi e^2\rho(\GG)/G^2$ & one FFT, \textbf{exact} \\
$\tauvW$ & $|\nabla\rho|^2/8\rho$ & one FFT, \textbf{exact} \\
\bottomrule
\end{tabular}
\end{center}

These are exact for band-limited fields --- precisely what a plane-wave grid
carries --- and the operator is already computing FFTs in every layer. The
physics and the architecture share one mathematical substrate. That is the
strongest technical argument for a Fourier neural operator here, independent of
accuracy.

\section{Roadmap}

\begin{center}
\begin{tabular}{@{}clc@{}}
\toprule
Phase & Work & Status \\
\midrule
0 & supervised baseline, honest split by material & done \\
1 & hard constraints: $\log$ head, charge normalisation, Pauli head & Pauli head done \\
2 & $H^1$ data loss; scaling-relation augmentation & available \\
3 & Euler--Lagrange residual with the analytic surrogate & available \\
4 & $\delta \Ts/\delta\rho$ by autograd; joint training & planned \\
5 & cycle consistency via density-to-potential inversion & planned \\
6 & deep-equilibrium orbital-free solver & stretch \\
\bottomrule
\end{tabular}
\end{center}

Phases~1--3 are low risk and independently valuable. Phase~4 is the research
contribution.
